\section{Problem formulation and experimental design}

\subsection{Scientific data and preprocessing}

% AUTHOR NOTE: Define the CAMELS data product, map dimensions, physical
% parameters, normalization, train/test partition, and the exact rule used to
% form every training subset. State explicitly that each real reference is
% computed from the subset used to train that model.
\draftplaceholder

\subsection{Unconditional novelty experiments}

% AUTHOR NOTE: Describe the UNet and DiT families, training-set sizes, optimizer
% update budgets, checkpoint selection, random seeds, and any controlled changes
% in width, depth, or patch size. Keep this separate from conditional calibration.
\draftplaceholder

\subsection{Conditional consistency experiments}

% AUTHOR NOTE: Define the full conditioning vector, requested parameter values,
% held-out simulation indices, generated samples per condition, and the frozen
% diagnostic model trained only on real fields. State how its held-out-real
% performance is validated before it is applied to generated fields.
\draftplaceholder

\subsection{Sampling and reproducibility controls}

% AUTHOR NOTE: Report the forward noise schedule, prediction parameterization,
% inference endpoint, DPM-Solver and DDPM controls, inference step counts, seeds,
% sample counts, and software or hardware details needed to reproduce results.
\draftplaceholder
